\section{Introduction}

The deployment of autonomous multi-agent systems into high-stakes environments has proceeded far ahead of our theoretical and architectural capacity to contain them. As these systems assume roles in domains spanning infrastructure management, financial operations, medical coordination, and legal reasoning, the consequences of uncontrolled agency cascades have shifted from theoretical concern to operational necessity. We argue that the absence of a mandatory, architecturally enforced bail-out mechanism represents a fundamental gap in the design of accountable autonomous systems, and we propose the \bxthree\ framework as a formal response to this gap.

\subsection{The Problem of Unbounded Agency}

Classical agent architectures treat autonomy as an open-ended property: an agent is provisioned with a goal, granted access to tools and environments, and permitted to pursue that goal with minimal external interruption. This design philosophy served early symbolic systems and remains appropriate for narrow, low-consequence applications. However, when agents operate in parallel, share environmental state, and mutate collective resources, the failure modes multiply in ways that single-agent frameworks fail to anticipate. A cascade in a multi-agent system is not merely the sum of individual failures; it is an emergent property of interaction dynamics that can exceed the fault tolerance of any individual component.

When systems lack a formal bail-out mechanism, uncontrolled cascades produce outcomes that are architecturally irreversible. The absence of a defined surrender point means that operators cannot intervene without engaging in ad hoc termination procedures that themselves introduce risk. We observe this pattern across existing deployments: shutdown protocols are implemented as afterthoughts, external to the agent's reasoning cycle, and therefore inaccessible to the agent itself at the moment when intervention is most needed.

\subsection{Accountability Without Architectural Fallback}

The notion of accountability in autonomous systems has been treated predominantly as a regulatory and ethical question. Regulatory frameworks ask who is responsible after the fact; ethical frameworks ask what the agent ought to have done. Neither framework supplies a mechanism that is operative within the agent's decision cycle. Accountability, to be effective, must be a structural property of the system, not a post hoc attribution. We contend that this requires an architectural fallback: a defined, enforced boundary that the system itself can recognize and respect.

The \bxthree\ framework introduces this fallback as a first-class construct. Where conventional systems separate safety constraints from the agent's core logic, \bxthree\ integrates the bail-out condition directly into the agent's operational semantics. This integration ensures that the fallback is not a patch applied by an external monitor but a property recognized by the agent as part of its own \purpose.

We identify four canonical failure modes that arise in the absence of such a fallback. First, \textbf{scope expansion}: agents pursue objectives beyond their defined \purpose\ because no internal mechanism recognizes the boundary. Second, \textbf{resource exhaustion}: agents consume computational, financial, or temporal resources without encountering a halting condition. Third, \textbf{accountability opacity}: when failure propagates across agent boundaries, no architectural trace exists to reconstruct causal responsibility. Fourth, \textbf{irreversibility}: systems cross into states from which recovery is impossible without external intervention that the architecture did not anticipate.

Each of these failure modes is directly mitigated by the presence of a mandatory, architecturally enforced bail-out condition. The condition operates as a \bounds\ primitive: it is not a soft constraint or a preference but a hard operational limit that overrides other directives when triggered.

\subsection{The \bxthree\ Framework}

The \bxthree\ acronym encodes the four primitives that constitute the framework. \purpose\ defines the categorical objective that constrains all agent activity; agents operating outside their \purpose\ are not merely suboptimal but architecturally unauthorized. \bounds\ defines the hard limits on resource consumption, temporal duration, and environmental interaction; when these limits are reached, the bail-out condition is triggered. \fact\ defines the evidentiary record that the system maintains; all agent actions generate \fact\ entries that together constitute a complete, auditable trail of decisions and their consequences. \bxthree\ itself is the integument that binds these primitives into a unified operational model.

The bail-out mechanism is not a shutdown command in the conventional sense. It is a structured transition in which the agent, upon reaching a \bounds\ condition, hands control to a designated fallback process that preserves \fact\ state, halts further action under the original \purpose, and reports to the operator with a complete accountability record. This transition is deterministic, repeatable, and verifiable. It does not require operator intervention to execute; it requires only operator notification afterward.

\subsection{Paper Roadmap}

The remainder of this paper is organized as follows. \S2 reviews related work in agent safety architectures, multi-agent failure modes, and accountability frameworks, distinguishing prior approaches from the \bxthree\ primitives. \S3 formalizes the \bxthree\ model, providing formal semantics for \purpose, \bounds, and \fact, and proving the conditions under which the bail-out transition is triggered. \S4 describes the implementation of the \bxthree\ runtime and evaluates its performance across a suite of multi-agent stress tests. \S5 discusses the implications of \bxthree\ for regulatory frameworks and the design of human-agent oversight structures. \S6 concludes with open questions and directions for future research.